%=====================================================================
\section{Introduction}
\label{sec:intro}

\begin{figure*}[t]
\centering
\resizebox{\textwidth}{!}{%
\begin{tikzpicture}[
    font=\small,
    >={Stealth[length=2mm]},
    rule/.style={draw, rounded corners=2pt, fill=gray!15, minimum height=6mm, inner sep=3pt, align=center, font=\footnotesize},
    dec/.style={draw, rounded corners=2pt, fill=blue!10, minimum height=6mm, inner sep=3pt, align=center, font=\footnotesize},
    box/.style={draw, rounded corners=3pt, inner sep=5pt, align=left},
    ttl/.style={font=\small\bfseries, anchor=south west},
  ]
  % ---------------- (a) the procedure
  \node[rule, text width=3.6cm] (cand) at (2.2,0) {\textbf{rules} enumerate candidates\\(verb + particle, light verb + noun, \dots)};
  \node[dec, text width=3.6cm] (ask) at (2.2,-1.35) {\textbf{closed questions}, one per candidate, options from the guideline};
  \node[rule, text width=3.6cm] (asm) at (2.2,-2.7) {\textbf{rules} read the probabilities and write the annotation};
  \draw[->] (cand) -- (ask); \draw[->] (ask) -- (asm);
  \node[font=\scriptsize, anchor=west, text width=2.3cm] at (4.25,-1.35) {Jev returns a distribution, never text};
  \node[ttl] at (-0.1,0.55) {(a) A guideline as a procedure};
  % ---------------- (b) a closed decision
  \node[box, fill=blue!4, text width=6.3cm, anchor=north west] (q) at (6.9,0.45) {%
    \emph{She finally \textbf{gave up} smoking after the operation.}\\[3pt]
    \textbf{Candidate} (rule): \emph{gave}\,(3) + \emph{up}\,(4), kind \texttt{particle}\\[2pt]
    \textbf{Q} What is the combination \emph{gave + up}?\\
    \quad VPC.full: \textbf{.86} \quad VPC.semi: .09 \quad none: .05\\[3pt]
    \textbf{Rule} keeps the best category if $p \geq \theta$: \texttt{3:VPC.full 4:VPC.full}.\\[2pt]
    \textbf{Cascade}: if $p < \tau$, the same question goes to a cheap LLM.};
  \node[ttl] at ([yshift=2pt]q.north west) {(b) A closed decision};
  % ---------------- (c) learning in batches
  \node[box, fill=green!6, text width=5.2cm, anchor=north west] (l) at (14.0,0.45) {%
    \textbf{Each round}, 5 training sentences:\\[2pt]
    1.~parse with the procedure;\\
    2.~the \textbf{reflector} reads gold vs.\ output and every question with its answers, and proposes $\leq$3 edits (an option, a question, a rule, a threshold);\\
    3.~each edit is \textbf{measured} on the batch, the two previous batches and a dev set, and kept only if it helps.\\[3pt]
    \textbf{Result}: the procedure, as a sequence of small tested diffs.};
  \node[ttl] at ([yshift=2pt]l.north west) {(c) Learning the procedure};
  \draw[->, thick] (ask.east) ++(2.3,0) -- ([yshift=-6mm]q.west);
  \draw[->, thick] (q.east |- l.west) -- (l.west);
  \coordinate (low) at ([yshift=-3mm]l.south);
  \draw[->, thick, dashed] (l.south) -- (low) -| node[font=\scriptsize\itshape, pos=0.25, below] {edits: the next version of the procedure} (asm.south);
\end{tikzpicture}}
\caption{Decision orchestration. (a)~A guideline is written as a procedure in which rules enumerate candidates and assemble the output, and a decision model answers closed questions with probabilities. (b)~One closed decision in the verbal-multiword-expression procedure. (c)~The procedure is learned from a few annotated sentences at a time: an LLM proposes edits, and each edit is kept only if it measurably helps.}
\label{fig:overview}
\end{figure*}

A new annotation scheme, or an old one in a new language or domain, usually starts with a written guideline and very little annotated data. Training a supervised model needs thousands of examples; prompting an LLM to write the annotation needs none, but hands the model the whole job: deciding what to annotate, applying the scheme's conventions, and producing well-formed output, all at once and at the price of a generative model for every sentence.

Most of a guideline is not judgment. It tells the annotator where to look, which units to consider, which tests to apply in which order, and how to write the result; only some steps require understanding the sentence. We propose \textbf{decision orchestration} (Figure~\ref{fig:overview}): write the guideline as an executable \emph{procedure} in which rules enumerate candidates and assemble the output, and a \emph{decision model} answers only closed questions, each a choice among options taken from the guideline, with a probability distribution. The output is well-formed by construction. Every decision is cheap, inspectable and comes with a confidence. Uncertain decisions can be routed to a stronger model.

A hand-written procedure encodes the guideline as its author read it. Gold annotations encode the conventions annotators actually followed, and the two differ. We therefore \emph{learn the procedure}, not model weights. An LLM reflector reads five training sentences at a time, with the gold and predicted annotations and every question asked with its answer, and proposes small edits to the procedure's code, options and thresholds. Each edit is kept only if it improves the score on the batch, on earlier batches and on a development set. The procedure's history becomes a sequence of small tested diffs that a person can read, and the whole learning process costs a few dollars.

We test the method on two schemes that differ in how much of the output is structure: verbal multiword expressions (PARSEME), where an annotation is a set of categorized token groups, and Uniform Meaning Representation (UMR), where it is a graph. Each is tested in Chinese and English. Our findings:
\begin{itemize}[leftmargin=*,itemsep=2pt,topsep=2pt]
\item On PARSEME, a procedure learned from \textbf{60 annotated sentences} reaches 46.3 F on English, \textbf{13 points above the best system of the 2018 shared task}, which trained on all 3,471 training sentences, and 57.8 F on Chinese, \textbf{18 points above} Claude Haiku annotating directly with ten gold examples. Twelve rounds of reflection add 5.5 F on Chinese (7.8 on unseen expressions) for \$0.53 of reflector calls, and the decisions for the whole Chinese test set cost \$0.06.
\item On UMR, where the output is mostly structure, closed decisions alone do not beat direct generation. Routing uncertain decisions adds 2.5 F in Chinese, and reflection adapts a procedure written for Chinese to English (+4.1 F); but direct generation with ten examples matches the best procedure in Chinese and exceeds it by 16 F in English.
\item The procedure's options must follow the annotation's conventions, not the full inventory of the guideline: a UMR procedure that offered the decision model the 27-value aspect lattice of the guidelines, where the gold uses seven coarse values, gained 5.2 F when its aspect options were restricted to those seven and one polarity rule was restored.
\end{itemize}
Decision orchestration thus pays off where a guideline is a sequence of tests over enumerable candidates, and it tells us, through the measured edits and the routed decisions, where a scheme is not.
